\documentclass[10pt,a4paper]{amsart}
\usepackage[margin=19mm]{geometry}
\usepackage{amsmath,amssymb,mathtools}
\usepackage[scr=rsfs,cal=euler]{mathalfa}
\usepackage{xfrac}
\usepackage[unicode,hidelinks,bookmarks=false]{hyperref}
\newcommand{\ie}{i.e.\ }
\newcommand{\cf}{cf.\ }
\newcommand{\resp}{respectively\ }
\newcommand{\Subsection}{\subsection}
\providecommand{\nobreakdash}{\nobreak\hbox{-}}
\setlength{\emergencystretch}{2em}
\multlinegap0pt
\usepackage{amssymb}
\usepackage{nicefrac}
\usepackage[all]{xy}
\usepackage{enumerate}
\usepackage{caption}
\usepackage{tikz}
\usepackage{subfig}
\newtheorem{proposition}{Proposition}
\newtheorem{lemma}{Lemma}
\newcommand{\NN}{{\mathbb N}}
\newcommand{\ZZ}{{\mathbb Z}}
\title{New constraints on Lagrangian embeddings, and the shape invariant}
\author{Richard Hind, Ely Kerman}
\date{Source: arXiv:2602.10291v1 (CC BY 4.0)}
\begin{document}
\maketitle

\noindent\emph{Source and licence.} New constraints on Lagrangian embeddings, and the shape invariant. Version arXiv:2602.10291v1 (CC BY 4.0), distributed by arXiv under the Creative Commons Attribution 4.0 International licence, http://creativecommons.org/licenses/by/4.0/, as recorded in the arXiv licence metadata for this exact version and shown on the abstract page https://arxiv.org/abs/2602.10291v1. Full licence notice: \texttt{licenses/arxiv-2602.10291-CC-BY-4.0.txt}.\par\smallskip

\noindent\textit{Editorial scope.} Counted unit: the article's lemma lem:pre (source0.tex lines 765--767). The article's complete original proof of it is delivered with it (source0.tex lines 769--799, one \texttt{proof} environment of 31 lines). No mathematics is written, repaired or re-derived: the statement and the proof are the article's own lines, in the article's own notation, and no retained line is substituted or re-flowed.  Retained uncounted as the source-local context the proof needs: lines 685--744.  Nothing was omitted from any retained span, so the package carries no omission record.  No retained line contains a citation, so the wrapper carries no bibliography and no bibliography entry.  No retained line contains a figure environment, an image-inclusion command or drawing code, so no image is delivered and the package declares no asset dependency.  Editorial formatting only, in addition: an AMS \texttt{amsart} preamble that restates the article's own declarations and macros used by the retained lines, namely the environments \texttt{proposition}, \texttt{lemma} on separate counters, together with the standard packages those lines need.  The retained environments print in this document as Proposition 1, Lemma 1 in order of appearance, whereas the article prints the same items with the numbering of its own full document.  The complete native source of the article as distributed in the arXiv e-print for this exact version is bundled unchanged as \texttt{sources/194381/source0.tex}.
\begin{proposition}\label{prop:submanifolds}
For every sufficiently large $d \in \NN$, there exist symplectic embeddings
\begin{align*}
    &F,G \colon S^2 \to X_2 \smallsetminus (L_1 \cup L_2) \\
    &D_A, D^0_2, D^\infty_2 \colon  (D^2,S^1) \to (X_2,L_2) \\
    &D_C, D^0_1, D^\infty_1 \colon  (D^2,S^1) \to (X_2,L_1) \\
    &O_B \colon (S^1 \times [0,1],S^1 \times \{0,1\}) \to (X_2,L_1 \cup L_2)
\end{align*}
such that:\\

\begin{enumerate}

\item  The maps $F$ and $G$ both represent the class $(d,1) \in H_2(X_2; \ZZ)$\\

\item The six disks all have Maslov index $2$.\\



\item The set $\{[\partial D_A], [\partial D^\infty_2]\}$ is an integral basis of $H_1(L_2;\ZZ)$, and $[\partial  D^0_2]=-[\partial  D^\infty_2]$.\\

\item The set $\{[\partial D_C], [ \partial D^\infty_1]\}$ is an integral basis of $H_1(L_1;\ZZ)$, and $[\partial  D^0_1]=-[\partial D^\infty_1]$.\\

\item Denoting the component of the boundary of $O_B$ on $L_i$ by $\partial^i O_B$, the set $\{[\partial^i O_B], [\partial D^\infty_i]\}$ is an integral basis of $H_1(L_i;\ZZ)$\\



\item The images of  $p \circ D_A$, $p \circ O_B$ and  $p \circ D_C$ are the closures of $A$, $B$ and $C$, respectively.\\

\item The symplectic areas of  $D_A$, $O_B$ and  $D_C$ are $s$, $r$ and $b-s$, respectively. The disks $D^\infty_1$ and $D^\infty_2$ both have symplectic area equal to $r$.\\

\item  All nine maps are $J_2$-holomorphic away from arbitrarily small neighborhoods of a collection of Lagrangian tori whose elements are close to, and Lagrangian isotopic to, either $L_1$ or $L_2$.\\

\item We have the following intersection pattern:\\
\begin{enumerate}
    \item $F \bullet G=2d$
    \item $F \bullet D^0_i + F \bullet D^\infty_i = G \bullet D^0_i + G \bullet D^\infty_i =1$ for $i=1,2$
    \item $\{F \bullet D^\infty_1, G \bullet D^\infty_1\} = \{F \bullet D^\infty_2, G \bullet D^\infty_2\} = \{0,1\}$    %$F \bullet D^\infty_1=1$ and $F \bullet D^\infty_2=0$
    \item $F \bullet D_A + G \bullet D_A = d$
    \item $F \bullet D_C + G \bullet D_C = d$
    \item $F \bullet O_B =1$ and $G \bullet O_B=0$
    \item $D^\infty_1 \bullet S_\infty = D^\infty_1 \bullet S_\infty =1$
    \item $F \bullet T_0 = F \bullet T_\infty=G \bullet T_0 = G \bullet T_\infty=1$.\\
\end{enumerate} 



% \item Exactly one of $F$ and $G$ intersect the planes of $\mathfrak{r}_0$ and the other intersects the planes of $\mathfrak{r}_\infty$.\\

% \item Exactly one of $F$ and $G$ intersect the planes of $\mathfrak{s}_0$ and the other intersects the planes of $\mathfrak{s}_\infty$.\\


% \item  $$F \bullet D_A + G \bullet D_A = d$$ $$F \bullet D_C + G \bullet D_C = d$$ $$F \bullet O_B =1$$ $$G \bullet O_B=0$$. 

\item The spheres $F$ and $G$ have exactly $2d$ intersection points. Of these, $d$ are located in $p^{-1}(A)$ and $d$ are located in $p^{-1}(C)$.\\

\end{enumerate} 



\end{proposition}

\begin{lemma}\label{lem:pre}
    The Lagrangian tori $L_1$ and $L_2$ are relatively exact in $(X_4, \Omega_4)$ and both are homologous to the zero section.
\end{lemma}

\begin{proof}

To prove that $L_1$ and $L_2$ are relatively exact, it suffices to find two relative cycles,  $C_1$ and $C_2$, representing classes in $H_2(X_4, L_1 \cup L_2;\ZZ)$, such that 
\begin{equation}
\Omega_4(C_1) = \Omega_4(C_2) =0
\end{equation}
\begin{equation}
\{[\partial C_1 \cap L_1], [\partial C_2 \cap L_1]\}\quad \text{is a basis of} \quad   H_1(L_1;\ZZ),
\end{equation}
and 
\begin{equation}
\{[\partial C_1 \cap L_2], [\partial C_2 \cap L_2]\} \quad \text{is a basis of} \quad H_1(L_2;\ZZ).
\end{equation}

To produce these cycles we will use the transforms of the submanifolds of $X_2$ described in Proposition \ref{prop:submanifolds}.

By Proposition \ref{prop:submanifolds} we can choose disks $D_1$ equal to one of the transforms $\hat{D}^0_1$ or $\hat{D}^\infty_1$, and $D_2$ equal to one of the transforms $\hat{D}^0_2$ or $\hat{D}^\infty_2$, such that both disks intersect $\hat{F}$ exactly once but are disjoint from $\hat{G}$. We denote the intersection points with $\hat{F}$ by $q_1$ and $q_2$ respectively.
%Consider the transforms $\hat{D}^\infty_1$ and $\hat{D}^\infty_2$ and let $q_1$ and $q_2$ be their respective intersection points 
%with $S_{\infty}$. 
Let $\gamma_1$ be a smooth curve in $\hat{F}$ connecting $q_1$ to $q_2$, and $\Gamma_1$ a smooth cylinder in the normal bundle to $\hat{F}$ such that $\Gamma_1$ is disjoint from $\hat{F}$ and projects to $\gamma_1$. We may assume $\Gamma_1$ intersects $D_1$ and $D_2$ in circles around their intersections with $\hat{F}$. The circles bound disks $E_1$ and $E_2$ such that $D_1 \smallsetminus E_1$ and $D_2 \smallsetminus E_2$ are cylinders disjoint from $\hat{F}$, and we can define $C_1$ to be the glued cycle $C_1 = (D_1 \smallsetminus E_1) \cup \Gamma_1 \cup (D_2 \smallsetminus E_2)$. %can then be perturbed to be disjoint from $\hat{F} \cup  \hat{G} \cup \hat{T}_0 \cup \hat{T}_{\infty}$ to produce our first cycle $C_1$. 
It clearly has a zero $\Omega_4$-area. Moreover,  $[\partial C_1 \cap L_1]=-[\partial D_1]$ and $[\partial C_1 \cap L_2]=[\partial D_2]$.

To construct the cycle $C_2$, we first consider the transform $\hat{O}_B$ of $O_B$. It follows from assertion (5) of Proposition \ref{prop:submanifolds} that 
$\{[\partial C_1 \cap L_i], [\partial \hat{O}_B \cap L_i]\}$ is a basis of  $H_1(L_i;\ZZ)$ for $i=1,2$.
However the $\Omega_4$-area of $\hat{O}_B$ is $r$. Moreover, $\hat{O}_B$ intersects the axis $\hat{F}$ positively in precisely one point and so does not represent a class in $H_2(X_4, L_1 \cup L_2;\ZZ)$. Both problems can be fixed simultaneously by choosing a smooth curve $\gamma_2$ connecting $q_1$ to $\hat{O}_B \cap \hat{F}$. We can then form a cylinder $\Gamma_2$ over $\gamma_2$ and use it to glue $D_1$ and $\hat{O}_B$ to produce the required cycle $C_2$ in the complement of $\hat{F}$ as above.
%replacing $\hat{O}_B$ with $(\hat{D}^\infty_1)^- \# \gamma_2 \# \hat{O}_B$ where $\gamma_2$ is a smooth curve in $\hat{F}$ from the point  $\hat{D}^\infty_1 \cap \hat{F}$ to the point $\hat{O}_B \cap \hat{F}$. 
%The resulting pair-of-pants map can then be perturbed off of $\hat{F}$ so that its image lies in $X_4$. This is the required cycle, $C_2$.


It remains to show that $L_1$ (or $L_2$) is homologous to the zero section. It is enough to show that the inclusion $H_1(L_1;\ZZ) \to H_1(X_4;\ZZ)$ is injective or, equivalently, that the boundary map $$\partial \colon H_1(X_4,L_1;\ZZ) \to H_1(L_1;\ZZ)$$ is trivial. If the boundary map is not trivial then there is a relative 2-cycle $C$ such that $[\partial C] = (m,n) \neq (0,0) \in H_1(L_1, \ZZ).$ Note that $C$ also represents a class in $H_2(X_3, L_1;\ZZ)$. Hence, adding suitable multiples of $\hat{D}_C$ and $\hat{D}_1^{\infty}$ to $C$ we get a closed cycle $C_{m,n}$ in $X_3$ whose intersections with the topological axes are determined by the disks added. If $m\neq 0$, then $C_{m,n}$, like $\hat{D}_C$,  will intersect $\hat{T}_0$ but not $\hat{T}_{\infty}$. Otherwise, if $m=0$ and $n\neq0$, then $C_{m,n}$, like $\hat{D}_1^{\infty}$,  will intersect $\hat{F}$ but not $\hat{G}$. In both cases, the resulting intersection pattern is impossible since $[\hat{T}_0]=[\hat{T}_{\infty}] \in H_2(X_3;\ZZ)$ and $[\hat{F}] =[\hat{G}]\in H_2(X_3;\ZZ)$.
\end{proof}

\end{document}
